
%% ===== 2026 规范 AI 工具使用规定第 3 条 =====
%% 「AI 工具使用声明」必须放在参考文献之前，按下面二选一模板写
%% (1) 未使用 AI 工具的参赛队，删除下面整段 \subsection*{AI 工具使用声明} + 段落
%%     并把第二段文字改为：
%%       \textbf{本参赛队在竞赛过程中未使用任何 AI 工具。}
%% (2) 使用 AI 工具的参赛队（默认），保持下面这段话：
%%       \textbf{本参赛队在竞赛过程中使用了 AI 工具，
%%         主要用于【语言润色、代码调试】，详细使用情况见支撑材料。}
\subsection*{AI 工具使用声明}

\textbf{本参赛队在竞赛过程中使用了 AI 工具，主要用于语言润色与代码调试，详细使用情况见支撑材料。}

%% 详细使用情况填到 \texttt{支撑材料/AI工具使用详情.pdf}（2026 AI 规定第 4 条），
%% 模板字段：
%%   (1) 工具名称、版本或型号
%%   (2) 具体使用目的和环节（按问题X/步骤Y分述）
%%   (3) 主要提示方式与使用过程（可附典型交互示例）
%%   (4) 对 AI 输出的采纳、人工修改和核验情况（语言润色除外）

\begin{thebibliography}{99}

\bibitem{wang2023jiyu} 王浩, 李明明. 基于梯度提升回归的零售销售额预测研究[J]. 统计与决策, 2023, 39(8): 75-80.

\bibitem{zhang2022tezheng} 张伟, 陈静. 特征工程在零售数据分析中的应用综述[J]. 计算机应用研究, 2022, 39(5): 1281-1288.

\bibitem{liu2023jicheng} 刘洋, 赵晓明. 基于集成学习的多因素销售预测方法[J]. 管理工程学报, 2023, 37(3): 134-142.

\bibitem{zhou2021shuju} 周明辉. 数据驱动的零售运营优化研究[D]. 北京: 清华大学, 2021.

\bibitem{chen2022minzu} 陈思远, 黄晓峰. 基于随机森林和Pearson相关的特征选择方法比较[J]. 数据分析与知识发现, 2022, 6(4): 52-63.

\bibitem{li2024jiyu} 李明, 王芳. 基于XGBoost的零售业销售预测模型研究[J]. 商业经济研究, 2024, (3): 156-159.

\bibitem{breiman2001random} Breiman L. Random forests[J]. Machine Learning, 2001, 45(1): 5-32.

\bibitem{friedman2001greedy} Friedman J H. Greedy function approximation: A gradient boosting machine[J]. Annals of Statistics, 2001, 29(5): 1189-1232.

\bibitem{behzadian2012topsis} Behzadian M, Otaghsara S K, Yazdani M, et al. A state-of-the-art survey of TOPSIS applications[J]. Expert Systems with Applications, 2012, 39(17): 13051-13069.

\bibitem{hastie2009elements} Hastie T, Tibshirani R, Friedman J. The Elements of Statistical Learning: Data Mining, Inference, and Prediction[M]. 2nd ed. New York: Springer, 2009.

\end{thebibliography}
